% !TeX root = ../../Book2.tex
\chapter{Hardy 空间、BMO 与加权调和分析}
\label{b2:app:C}

\begin{chapterabstract}

强型 \(L^1\) 与 \(L^\infty\) 端点失效后，原子 Hardy 空间和 BMO 提供保留抵消的替代空间。本附录证明 John–Nirenberg 指数尾、原子空间的 BMO 对偶与奇异积分的原子界，并介绍极大截断、\(T1\) 思想及 \(A_p\) 权的定义和必要性。

\end{chapterabstract}

\begin{learninggoals}\begin{itemize}
\item\label{b2:item:app-C-goal-1}用极大截断的振荡估计证明主值存在性的稠密性传递。
\item\label{b2:item:app-C-goal-2}证明 BMO 的指数尾并解释模常数的对偶对象。
\item\label{b2:item:app-C-goal-3}构造原子级数空间，验证 BMO 配对和奇异积分延拓。
\item\label{b2:item:app-C-goal-4}检验权的局部可积性、\(A_p\) 必要条件与幂权范围，并证明 \(A_1\) 子类的充分界。
\end{itemize}\end{learninggoals}


本附录沿用\chapref{b2:ch:17}和\chapref{b2:ch:18}的极大函数与奇异积分约定。为使端点论证自足，先以平方可积原子定义 Hardy 空间；它与各种极大函数定义的等价是另一个深层结构问题，不在证明过程中预设。

\section{极大截断与端点空间}

\label{b2:sec:app-C-1}


\chapref{b2:ch:18}的极大截断 \(T^*f=\sup_{\epsilon>0}|T_\epsilon f|\) 控制整个主值族。若截断在 \(L^2\) 中实现给定算子，且 \(T^*\) 在某个 \(L^p\)、\(p>1\) 上有界，便可从光滑稠密输入的主值存在性推广到全部 \(L^p\)。关键是
\[
 \limsup_{\epsilon,\delta\downarrow0}|T_\epsilon f-T_\delta f|
 \le2T^*(f-g)+
   \limsup_{\epsilon,\delta\downarrow0}|T_\epsilon g-T_\delta g|.
\]
先选光滑 \(g\) 逼近，再估计右侧异常集合。这是振荡与稠密性结合的证明，不能仅由各个截断分别有界推出全参数极限。

\ProseParagraphBreak
对于端点 \(L^1\)，主值与弱型理论仍可讨论，但强型范数不成立。原子空间在输入中加入局部零平均，BMO 则允许输出有不可避免的对数振荡。它们保留的是抵消，而不是简单地改变一个可积指数。


\section{BMO 与指数尾}

\label{b2:sec:app-C-2}


\begin{definition}\label{b2:def:C-bmo}
对 \(b\in L^1_{\mathrm{loc}}(\mathbb R^d)\)，定义
\[
 \|b\|_{\mathrm{BMO}}=\sup_Q\frac1{|Q|}\int_Q|b-b_Q|,
 \qquad b_Q=\frac1{|Q|}\int_Q b,
\]
上确界取所有轴平行立方体。有限者构成
\term[youjie-pingjun-zhendang]{有界平均振荡空间}[bounded mean oscillation]，常数函数视为同一个零类。
\end{definition}
\symidx{\(\mathrm{BMO}\)：有界平均振荡空间，按模常数解释}
这是半范数；若半范数为零，每个立方体上函数等于一个常数，重叠立方体使这些常数相同，故函数几乎处处为常数。任意有界函数有 BMO 半范数不超过 \(2\|b\|_\infty\)，但反向不成立。

下面采用\term[john-nirenberg-budengshi]{John--Nirenberg 不等式}[John--Nirenberg inequality]的下述版本。

\begin{theorem}[John--Nirenberg]\label{b2:thm:C-john-nirenberg}
存在只依赖维数的 \(c,C>0\)，使
\[
 |\{x\in Q:|b(x)-b_Q|>t\}|
 \le C|Q|e^{-ct/\|b\|_{\mathrm{BMO}}}.
\]
特别地，对每个 \(1\le r<\infty\)，
\[
 \left(\frac1{|Q|}\int_Q|b-b_Q|^r\right)^{1/r}
 \le C_{d,r}\|b\|_{\mathrm{BMO}}.
\]
\end{theorem}
\begin{proof}
归一化半范数为1。在 \(Q\) 内，对 \(|b-b_Q|\) 作二进极大立方体分解，阈值取 \(\lambda>2\)。坏立方体总测度至多 \(|Q|/\lambda\)，好集上函数绝对值至多 \(\lambda\)，且每个坏块 \(Q_j\) 的平均绝对值不超过 \(2^d\lambda\)，所以
\(|b_{Q_j}-b_Q|\le2^d\lambda\)。
在每个坏块上对 \(b-b_{Q_j}\) 重复同样操作。第 \(k\) 代坏集测度至多 \(\lambda^{-k}|Q|\)；离开这代坏集的点，通过逐代平均值差的望远镜相加，满足
\(|b-b_Q|\le k2^d\lambda\)（调整初始常数即可）。
因此在 \(t\) 的相邻整数代之间选 \(k\)，得到指数尾。零半范数情形已单独说明。
最后用层蛋糕
\(\int_Q|b-b_Q|^r=r\int_0^\infty t^{r-1}|\{|b-b_Q|>t\}|\,\mathrm dt\)，指数积分有限，便得到所有有限阶矩。
\end{proof}
函数 \(\log|x|\) 是典型无界 BMO 例子。若立方体远离原点，梯度 \(|x|^{-1}\) 与边长的乘积有界；若立方体靠近原点，按边长缩放，把常数 \(\log\ell(Q)\) 减去后，问题落在固定大小、位于固定有界区域的立方体上。那里 \(|\log|x||\) 可积，平均振荡统一有界。这个函数的指数尾与定理相符；它不属于 \(L^\infty\)。


\section{原子 Hardy 空间及其对偶}

\label{b2:sec:app-C-3}


\begin{definition}\label{b2:def:C-atom}
一个\term[pingfang-keji-yuanzi]{平方可积原子}[L2 atom] \(a\) 支撑于某立方体 \(Q\)，满足
\(\int a=0\)、\(\|a\|_2\le|Q|^{-1/2}\)。
对在 \(L^1\) 中可写成绝对可求和原子级数的函数定义范数
\(\|f\|_{\mathrm{at}}=\inf\sum|\lambda_j|\)，其中
\(f=\sum\lambda_ja_j\)、\(\sum|\lambda_j|<\infty\)；这些函数组成\term[yuanzi-hardy-kongjian]{原子 Hardy 空间}[atomic Hardy space]，记为 \(H^1_{\mathrm{at}}\)。
\end{definition}
\symidx{\(H^1_{\mathrm{at}}\)：原子 Hardy 空间}
每个原子有 \(\|a\|_1\le1\)，所以原子范数控制 \(L^1\) 范数，且每个元素积分为零。这个空间完备：对原子范数 Cauchy 列，取相邻差范数之和有限的快速子列，每个差选择系数和不超过其范数加一个可求和误差的表示。把首项与全部差的级数合并，得到一个绝对可求和原子表示；它同时给出 \(L^1\) 极限，并使余项原子范数趋零。有限原子和在此空间稠密，因为每个表示的系数尾和趋零。范数为零蕴含 \(L^1\) 范数为零，三角不等式由合并级数得到。

\begin{theorem}[原子空间的 BMO 对偶]\label{b2:thm:C-atomic-duality}
通过原子和上的配对 \(\Lambda_b(f)=\int fb\)，
\((H^1_{\mathrm{at}})^*\) 与 \(\mathrm{BMO}/\{\text{常数}\}\) 等价同构，范数相互控制。
\end{theorem}
\begin{proof}
John–Nirenberg 的 \(r=2\) 结论和零平均给出
\[
 |\int ab|=|\int_Q a(b-b_Q)|
 \le\|a\|_2\|b-b_Q\|_2\le C\|b\|_{\mathrm{BMO}}.
\]
对原子级数，定义 \(\Lambda_b(f)=\sum_j\lambda_j\int a_jb\)，上界保证绝对收敛。还需检查表示无关性。先对实 \(b\) 作有界截断 \(b_N=\max\{-N,\min\{b,N\}\}\)。截断为一 Lipschitz 函数，故 \(\|b_N\|_{\mathrm{BMO}}\leq2\|b\|_{\mathrm{BMO}}\)：以截断后的 \(b_Q\) 作比较常数，再换成真正平均至多再损失因子二。对复 \(b\) 分别截断实部、虚部，常数仍统一。若一个原子级数在 \(L^1\) 中和为零，则对有界 \(b_N\) 的配对为零；逐个原子用局部 \(L^2\) 收敛，全部级数用统一原子界与可和系数，令 \(N\to\infty\)，得到其对 \(b\) 的配对也为零。因此定义不依赖表示。再取系数和的下确界，得到连续泛函。一般配对如此定义，不要求任意 \(fb\) 的绝对积分预先有意义。

反向设 \(\Lambda\) 为连续泛函。对支撑于 \(Q\)、零平均的 \(f\in L^2\)，把 \(f\) 除以 \(|Q|^{1/2}\|f\|_2\) 就是原子，所以
\(|\Lambda(f)|\le\|\Lambda\||Q|^{1/2}\|f\|_2\)。
在 \(L^2_0(Q)\) 上的 Riesz 表示给出一个局部 \(L^2\) 函数 \(b_Q\)，使 \(\Lambda(f)=\int fb_Q\)，并有
\(\|b_Q-(b_Q)_Q\|_2\le\|\Lambda\||Q|^{1/2}\)。
对包含关系 \(Q\subset Q'\)，两表示在 \(L^2_0(Q)\) 上相同，所以其差在 \(Q\) 上为常数。选一列递增立方体覆盖全空间，逐个调整常数，就拼成一个 \(b\in L^2_{\mathrm{loc}}\)。
对任意立方体，用同一 Riesz 界和 Cauchy–Schwarz 得
\(|Q|^{-1}\int_Q|b-b_Q|\le\|\Lambda\|\)，故 \(b\in\mathrm{BMO}\)。配对在所有原子上相同，因此在整个空间相同。若两个 \(b\) 表示同一泛函，其差在每个立方体与所有零平均 \(L^2\) 函数正交，故为全局常数。
\end{proof}
经典实变量 \(H^1\) 也可由 Poisson 极大函数或 Riesz 变换定义。其与上述原子空间的等价需要原子分解定理；本附录将它作为进一步阅读方向，见 \cite{Stein1993Harmonic,Grafakos2014Modern}，而不把该等价偷偷用在已证对偶中。


\section{奇异积分原子界与 \(T1\) 思想}

\label{b2:sec:app-C-4}


\begin{proposition}\label{b2:prop:C-atom-cz}
\chapref{b2:ch:18}的标准 Calderón–Zygmund 算子满足
\(\|Ta\|_1\le C\)，其中 \(a\) 为任意上述原子。故
\(T:H^1_{\mathrm{at}}\to L^1\) 有界。
\end{proposition}
\begin{proof}
在固定扩大立方体 \(Q^*\) 内，
\(\int_{Q^*}|Ta|\le|Q^*|^{1/2}\|T\|_{2\to2}\|a\|_2\le C\)。
在外部用原子零平均减去中心核：
\[
 Ta(x)=\int_Q[K(x,y)-K(x,c_Q)]a(y)\,\mathrm dy.
\]
核光滑条件积分后给出
\(\int_{(Q^*)^c}|Ta|\le C\|a\|_1\le C\)。
对有限原子和得范数界。为了确认绝对原子级数的输出与表示无关，可用 \(L^2\) 伴随测试：对 \(\varphi\in C_c^\infty\)，原子界给出 \(|\int a\,T^*\varphi|\leq C\|\varphi\|_\infty\)。具体地，在任意 \(Q\) 上令 \(b=T^*\varphi-(T^*\varphi)_Q\)，选原子 \(a=\overline b/(|Q|^{1/2}\|b\|_2)\)（零函数另处理），就得到 \(|Q|^{-1/2}\|b\|_2\leq C\|\varphi\|_\infty\)。Cauchy--Schwarz 因而直接给出 \(T^*\varphi\) 的 BMO 界。此处不预设一个已在整个原子空间上定义的泛函。若输入原子级数在 \(L^1\) 中和为零，其与该 BMO 函数的配对为零，故输出级数与所有测试函数配对为零。输出在 \(L^1\) 中绝对收敛，因此等于零。这建立表示无关性和有界延拓；在有限原子和上它与给定 \(L^2\) 算子相同。
\end{proof}
对偶路线提示 \(L^\infty\) 输出应由 BMO 控制，但要精确定义 \(T\) 对非衰减输入的作用，通常须减去远场常数，并按模常数解释。这一规范化不是普通绝对积分。

\ProseParagraphBreak
\(T1\) 思想反过来问：能否由核的标准条件、局部弱有界性以及 \(T1,T^*1\) 的 BMO 行为推出 \(L^2\) 有界性？其中常数输入检测低频抵消，弱有界性检测对角附近的局部作用，核条件检测支撑分离后的行为。三个部分相互补充。一般 \(T1\) 定理的完整算子演算不在此建立，系统理论可参看 \cite{Stein1993Harmonic}；不能只验证 \(T1=0\) 就认定所有必要假设已满足。


\section{权、\(A_p\) 条件与加权理论}

\label{b2:sec:app-C-5}


权 \(w\) 指几乎处处正、局部可积的函数，
\(\|f\|_{L^p(w)}^p=\int|f|^pw\)。对 \(1<p<\infty\)，定义
\[
 [w]_{A_p}=\sup_Q
 \left(\frac1{|Q|}\int_Qw\right)
 \left(\frac1{|Q|}\int_Qw^{-1/(p-1)}\right)^{p-1}<\infty.
\]
\symidx{\([w]_{A_p}\)：Muckenhoupt 权常数}
这称为\term[muckenhoupt-quan]{Muckenhoupt 权}[Muckenhoupt weight]条件。
\(A_1\) 则定义为 \(Mw\le Cw\) 几乎处处（使用一致的立方体极大算子）。

\begin{proposition}\label{b2:prop:C-ap-necessary}
若非中心极大算子在 \(L^p(w)\) 上有界，\(1<p<\infty\)，则 \(w\in A_p\)。
\end{proposition}
\begin{proof}
在立方体 \(Q\) 上取 \(f=w^{-1/(p-1)}\mathbf1_Q\)，先取 \(f_N=\min(w^{-1/(p-1)},N)\mathbf1_Q\)。它满足 \(f_N^{p-1}w\leq1\)，所以 \(\int f_N^pw\leq\int f_N\)；在有限 \(Q\) 上两者均有限。对 \(x\in Q\)，有 \(Mf(x)\ge |Q|^{-1}\int_Qf\)。
代入范数界得到
\[
 |Q|^{-p}\left(\int_Qw^{-1/(p-1)}\right)^p\int_Qw
 \le C^p\int_Qw^{-1/(p-1)}.
\]
截断极限用单调收敛，先确认该逆权积分有限：令 \(F_N=\int_Qf_N\)，截断界给出 \(|Q|^{-p}F_N^p w(Q)\leq C^pF_N\)。因 \(w(Q)>0\)，除去 \(F_N>0\) 后可得 \(F_N\) 的统一上界。除去正因子即得到 \(A_p\) 条件。
\end{proof}
对于幂权 \(w(x)=|x|^\alpha\)，条件恰为
\(-d<\alpha<d(p-1)\)。
必要性来自原点附近权与逆权的局部可积性。充分性分两类立方体：远离原点者上 \(|x|\) 彼此可比，两个平均的幂相消；靠近原点者包含于半径 \(C\ell(Q)\) 的球，积分分别至多
\(C\ell(Q)^{d+\alpha}\) 与
\(C\ell(Q)^{d-\alpha/(p-1)}\)，除以体积后乘积也有统一界。

\ProseParagraphBreak
一个更强但易于核验的权条件，可以直接给出充分性。设 \(w>0\) 几乎处，局部可积，并且对全部球 \(B\) 有
\[
 \frac1{|B|}\int_Bw\leq A\,\operatorname*{ess\,inf}_B w.
\]
称它满足球形式的 \(A_1\) 条件。球与立方体的固定倍包含关系使其与立方体定义等价，常数只差维数因子。对立方体形式，直接以本性下确界估计逆权平均，可得 \([w]_{A_p}\leq[w]_{A_1}\)、\(p>1\)。所以这是全部开区间权条件的一种充分子类。
\begin{proposition}\label{b2:prop:C-a1-maximal}
若 \(w\) 满足上述条件，则中心极大函数满足
\[
 \int_{\{Mf>\lambda\}}w\leq \frac{C_dA}{\lambda}\int|f|w.
\]
故对每个 \(1<p<\infty\)，\(M\) 在 \(L^p(w)\) 上有界。
\end{proposition}
\begin{proof}
先取紧支集输入，且 \(\int|f|w<\infty\)。每个有限球的权平均为正，\(A_1\) 因而给出该球上正的本性下界，所以输入局部无权可积。对超水平集合的紧子集取平均超过 \(\lambda\) 的有限球覆盖，选出两两不交的球 \(B_j\)，其五倍球覆盖原紧集。权条件给出
\[
 \begin{aligned}
 w(5B_j)
 &\leq A\,|5B_j|\operatorname*{ess\,inf}_{5B_j}w\\
 &\leq A\,5^d|B_j|\operatorname*{ess\,inf}_{B_j}w
 \leq \frac{A\,5^d}{\lambda}\int_{B_j}|f|w.
 \end{aligned}
\]
相加时不交性使右端不重复计数。用局部有限测度 \(w\,\mathrm dx\) 的正则性取紧集极限，再作输入及位置截断，就得到完整弱型界。无穷端点仍是平均不超过本性上确界；在测度 \(w\,\mathrm dx\) 上应用 Marcinkiewicz 即得强型。
\end{proof}
这一证明只处理 \(A_1\) 子类，没有完成全部 \(A_p\) 的充分性。一般权需要更精细的自改进或层集论证；标准奇异积分的全部加权界也还需要对应算子工具。进一步阅读见 \cite{Grafakos2014Modern,Duoandikoetxea2001Fourier}。任意正权都给出一个 \(L^p(w)\) 空间，却远非任意正权都允许原来的算子界；幂权和 \(A_1\) 条件提供了两种可以实际核验的检验。


\begin{chaptersummary}
极大截断把范数估计转化为逐点主值信息；BMO 记录局部平均振荡，原子的零平均使它与端点 Hardy 空间形成对偶。指数尾、原子配对的表示无关性及算子延拓均已证明。加权部分建立了 \(A_p\) 必要条件、幂权范围及 \(A_1\) 子类的充分界；全部加权充分性与极大函数原子等价仍是进一步理论。
\end{chaptersummary}

\BookExercises

% 自拟：检查本附录的已建工具及其假设边界。

\begin{exercise}\label{b2:ex:app-C-1}
设 \(a\) 为平方可积原子，\(y\in\mathbb R^d\)、\(r>0\)。证明
\(a_{y,r}(x)=r^{-d}a((x-y)/r)\) 仍是原子，并推出该变换在 \(H^1_{\mathrm{at}}\) 上等距。
\end{exercise}
\begin{solution}
新支撑为 \(y+rQ\)，积分仍为零，且
\(\|a_{y,r}\|_2=r^{-d/2}\|a\|_2\leq|rQ|^{-1/2}\)。
对原子级数逐项变换，系数不变、\(L^1\) 收敛不变，故新原子范数不超过原范数。逆变换的参数为 \((-y/r,1/r)\)，同样给出反向不等式，范数因而相等。振幅 \(r^{-d}\) 保持积分规模；若改用平方归一化 \(r^{-d/2}\)，就不再是这一等距变换。
\end{solution}

% 自拟：检查本附录的已建工具及其假设边界。

\begin{exercise}\label{b2:ex:app-C-2}
记 \(B=\|b\|_{\mathrm{BMO}}>0\)。若 John--Nirenberg 尾界为
\(|\{|b-b_Q|>t\}|\leq C|Q|e^{-ct/B}\)，证明对 \(0<\epsilon<c\)，
\[
 \frac1{|Q|}\int_Q e^{\epsilon|b-b_Q|/B}
 \leq1+\frac{C\epsilon}{c-\epsilon}.
\]
\end{exercise}
\begin{solution}
对 \(u\geq0\)，有 \(e^{\epsilon u}=1+\epsilon\int_0^u e^{\epsilon s}\,\mathrm ds\)。
将 \(u=|b-b_Q|/B\) 代入，以 Tonelli 换序，再用题设尾界，得到
\(1+C\epsilon\int_0^\infty e^{-(c-\epsilon)s}\,\mathrm ds\)，即所需常数。严格条件 \(\epsilon<c\) 使最后积分收敛；仅有所有有限阶矩可积，并不能省略指数参数的检查。
\end{solution}

% 自拟：检查本附录的已建工具及其假设边界。

\begin{exercise}\label{b2:ex:app-C-3}
正文已证明 \(\log|x|\in\mathrm{BMO}(\mathbb R^d)\)。证明其平方
\((\log|x|)^2\) 不在 BMO 中，说明 BMO 不在普通乘法下封闭。
\end{exercise}
\begin{solution}
固定远离原点的立方体 \(Q=(1,2)^d\)，令 \(\ell(z)=\log|z|\)。它在 \(Q\) 上非常数，故
\(A=|Q|^{-1}\int_Q|\ell-\ell_Q|>0\)，并令
\(D=|Q|^{-1}\int_Q|\ell^2-(\ell^2)_Q|<\infty\)。
在 \(rQ\) 上换元 \(x=rz\)，平方函数减去平均后成为
\[
 2\log r\,(\ell-\ell_Q)+\ell^2-(\ell^2)_Q.
\]
反三角不等式使平均绝对振荡至少为 \(2|\log r|A-D\)，随 \(r\downarrow0\) 趋于无穷。因此不存在统一 BMO 半范数。常数项 \((\log r)^2\) 已被平均消去，真正增长的是与非恒定局部振荡相乘的一阶项。
\end{solution}

% 自拟：检查本附录的已建工具及其假设边界。

\begin{exercise}\label{b2:ex:app-C-4}
在直线上令
\[
 f(x)=\frac{\mathbf1_{(0,e^{-2})}(x)}
 {x(\log(1/x))^2}-\frac12\mathbf1_{(1,2)}(x).
\]
证明 \(f\in L^1\)、\(\int f=0\)，但 \(f\notin H^1_{\mathrm{at}}\)。可使用正文的 BMO 对偶与有界截断。
\end{exercise}
\begin{solution}
换元 \(s=\log(1/x)\)，第一项积分为 \(\int_2^\infty s^{-2}\,\mathrm ds=1/2\)，故两个断言成立。取 \(b(x)=\log|x|\in\mathrm{BMO}\)，令 \(b_N\) 为它在 \([-N,N]\) 上的截断。正文给出 \(b_N\) 的 BMO 半范数一致有界。若 \(f\in H^1_{\mathrm{at}}\)，有界配对等于普通积分，故对偶界使 \(\int f b_N\) 一致有界。
然而对 \(N>2\)，正奇性部分的配对为
\[
 -\int_2^\infty\frac{\min(s,N)}{s^2}\,\mathrm ds
 =-\log(N/2)-1,
\]
另一段的积分固定有限，故总配对趋于负无穷，矛盾。零积分只是原子空间的必要条件，不能替代局部抵消的定量控制。
\end{solution}

% 自拟：检查本附录的已建工具及其假设边界。

\begin{exercise}\label{b2:ex:app-C-5}
证明幂权 \(w(x)=|x|^\alpha\) 属于 \(A_1\) 当且仅当
\(-d<\alpha\leq0\)，并说明这一结果为何可以产生全部 \(p>1\) 上的加权极大函数界。
\end{exercise}
\begin{solution}
权须局部可积，所以 \(\alpha>-d\)。若 \(\alpha>0\)，含原点的球中本性下确界为零，平均却为正，违反 \(A_1\) 球条件。若 \(\alpha=0\)，常数权显然满足条件。
对于 \(-d<\alpha<0\)，球心距原点至少为两倍半径时，球内距离可比，平均不超过本性下确界的固定倍数。其余球包含于半径至多三倍原半径的原点球，平均至多 \(Cr^\alpha\)；而球内距离至多 \(3r\)，因幂为负，本性下确界至少为 \((3r)^\alpha\)。两类都满足统一条件。正文的 \(A_1\) 弱型与插值随后给出每个 \(p>1\) 的强型界；常数允许依赖 \(\alpha,d,p\)。
\end{solution}

% 自拟：检查本附录的已建工具及其假设边界。

\begin{exercise}\label{b2:ex:app-C-6}
若 \(1<p<q<\infty\)，证明 \(A_p\subset A_q\)，且
\([w]_{A_q}\leq[w]_{A_p}\)。
\end{exercise}
\begin{solution}
置 \(\theta=(p-1)/(q-1)\in(0,1)\)、\(v=w^{-1/(p-1)}\)。
对立方体上的归一化测度，Hölder 或凹函数 Jensen 给出
\[
 \frac1{|Q|}\int_Qv^\theta
 \leq\left(\frac1{|Q|}\int_Qv\right)^\theta.
\]
将它提高到 \(q-1\) 次幂并乘 \(w\) 的平均，就被 \(A_p\) 表达控制；对 \(Q\) 取上确界得到结论。这是增大输入指数方向的包含，不允许由 \(A_q\) 反推较小的 \(A_p\)。幂权范围可直接给出反向不成立的例子。
\end{solution}

% 自拟：检查本附录的已建工具及其假设边界。

\begin{exercise}\label{b2:ex:app-C-7}

设 \(w\in A_p\)，证明 \(w^{-1/(p-1)}\in A_{p'}\)，并计算特征常数。

\end{exercise}

\begin{solution}

记 \(v=w^{-1/(p-1)}\)，因 \(p'-1=1/(p-1)\)，有 \(v^{-1/(p'-1)}=w\)。交换两个平均后，\(A_{p'}\) 表达为原 \(A_p\) 表达的 \(1/(p-1)\) 次幂，故 \([v]_{A_{p'}}=[w]_{A_p}^{1/(p-1)}\)。

\end{solution}

% 自拟：检查本附录的已建工具及其假设边界。

\begin{exercise}\label{b2:ex:app-C-8}

为什么 \(T1=0\) 不能代替奇异积分的 \(L^2\) 起点？

\end{exercise}

\begin{solution}

这个等式只检测一个常数输入的抵消，既未控制一般 \(L^2\) 输入，也未控制对角局部作用和伴随。一般 \(T1\) 判据还含标准核、弱有界性、\(T^*1\) 的BMO条件及作用定义。缺少这些条件不能应用该判据。

\end{solution}
